\documentclass[12pt]{article}
\usepackage[utf8]{inputenc}
\usepackage[T1]{fontenc}
\usepackage[margin=1in]{geometry}
\usepackage{setspace}
\usepackage{xcolor}
\usepackage[colorlinks=true, urlcolor=blue!60!black, linkcolor=blue!60!black]{hyperref}
\setlength{\parindent}{0em}
\setlength{\parskip}{0.6em}
\renewcommand{\baselinestretch}{1.0}
\begin{document}
\begin{center}
{\LARGE \textit{Luther v. Borden}, 48 U.S. (7 How.) 1 (1849)\par}
\vspace{0.5em}
{\large Case Brief --- Two-Page Edition\par}
\end{center}

\section{Memory Jogger}
Which of two rival state governments is lawful is a political question for Congress and the President, not the courts (\href{https://supreme.justia.com/cases/federal/us/48/1/}{48 U.S.\ 1}).

\section{Detailed Case Facts}
Rhode Island still ran under its 1663 charter, which limited voting to freeholders. Reformers adopted an extralegal ``People's Constitution'' and in 1842 elected Thomas Dorr governor. The charter government declared martial law and won President Tyler's recognition. Militiaman Luther Borden and others broke into Martin Luther's home to arrest him for aiding Dorr.

\section{Procedural History}
Luther sued for trespass in federal circuit court. Justice Story excluded his evidence that the People's Constitution was lawfully ratified, and the jury found for the defendants.

\section{Judicial Votes}
Affirmed, 5--1. Taney wrote for McLean, Wayne, Nelson, and Grier; Woodbury dissented in part; Catron, Daniel, and McKinley did not participate.

\section{Holding}
Federal courts cannot decide which regime is a state's lawful government or whether it is republican under the Guarantee Clause; that is for Congress and the President, whose recognition binds the courts. A state may declare martial law to suppress an insurrection, and officers may enter homes on reasonable grounds.

\section{Analysis of Opinions}
Taney warned that voiding the charter government would nullify a year of laws and judgments; noted that Rhode Island's courts had upheld it; found no manageable standards, since courts cannot count voters and juries could disagree; and held that the Guarantee Clause commits the question to Congress and to the President under the 1795 Act, whose decision was conclusive under \textit{Martin v.\ Mott}. Woodbury agreed on the political question but argued that no legislature could subject civilians to martial law while courts were open, rejecting advance sanction for ``the reign of the bayonet.''

\section{Examples: Future Applications}
\textbf{Same-Side:} a Guarantee Clause challenge to ballot initiatives is nonjusticiable (\textit{Pacific States}). \textbf{Opposite-Side:} Equal Protection malapportionment claims are justiciable (\textit{Baker v.\ Carr}). \textbf{Fence-Sitter:} a Guarantee Clause challenge to federal commandeering, left open in \textit{New York v.\ United States}.

\section{Critique}
Originalists such as Ryan Williams and Tara Leigh Grove question whether the Clause or early practice supports a jurisdictional bar. Erwin Chemerinsky argues \textit{Luther} erased the Guarantee Clause; Woodbury's dissent is praised as a defense of civil liberties.

\section{Key Quotations}
\begin{quote}
``Under this article of the Constitution, it rests with Congress to decide what government is the established one in a State. For as the United States guarantee to each State a republican government, Congress must necessarily decide what government is established in the State before it can determine whether it is republican or not.''
\end{quote}
\noindent --- Chief Justice Taney, 48 U.S.\ (7 How.)\ at 42.

\end{document}
